\begin{frame}{$\kappa \approx 46$ vs $\kappa \approx 2$: 같은 골짜기가
아님}
  \begin{center}
  \small
  \begin{tabular}{@{}lll@{}}
    \toprule
     & $x=[1,2,3]$ & ``균형 잡힌'' 데이터
      (100개, $x \sim \mathcal{N}(0,1)$) \\
    \midrule
    조건수 $\kappa$ & $\approx 46.1$ & $\approx 2$ \\
    등고선 & 긴 축:짧은 축 $\approx 46{:}1$ 타원 --- \kb{zigzag}하며
      천천히 하강 & 거의 원형 --- 직선에 가까운 경로 \\
    \bottomrule
  \end{tabular}
  \end{center}
  \vskip0.8em
  \begin{block}{스케일 10배가 조건수를 $\sim$71배로 만들면}
    같은 데이터에 특징 스케일 10배 ($x = [10,20,30]$, 참값
    $y = 0.2x + 1$ --- 직선은 \textbf{같음}) 로 바꾸면
    $\kappa: 46.1 \to 3278.7$ --- 약 \textbf{71배} 커짐.
    (bias 열은 스케일링되지 않아 정확히 100배는 아니고, 두 고유값이
    스케일에 따라 \emph{다르게} 늘어나 비율이 바뀜.)
  \end{block}
\end{frame}
